\makeatletter
\def\rap@fontpath{../../../Common/}
\makeatother
\documentclass{../../../Common/Latex/Templates/rapport}
% ---------------------------------------------------------------------------
% Shared settings for the H2026 General Assembly document set.
%
% Usage, directly after \documentclass{../../../Common/Latex/Templates/rapport}
% in a document one folder below this file:
%   % ---------------------------------------------------------------------------
% Shared settings for the H2026 General Assembly document set.
%
% Usage, directly after \documentclass{../../../Common/Latex/Templates/rapport}
% in a document one folder below this file:
%   % ---------------------------------------------------------------------------
% Shared settings for the H2026 General Assembly document set.
%
% Usage, directly after \documentclass{../../../Common/Latex/Templates/rapport}
% in a document one folder below this file:
%   \input{../ga-common}
%
% Image paths are relative to the compiled document's folder, so every
% document using this file must sit one folder below it.
% ---------------------------------------------------------------------------

% Meeting details, used on covers and in each document's ID block.
\newcommand{\meetingdate}{30 September 2026}
\newcommand{\meetingtime}{12:00}
\newcommand{\meetinglocation}{Festsalen, NMBU}
\newcommand{\meetingline}{\textbf{Meeting:} \meetingdate, \meetingtime, \meetinglocation}

% Header: moon logo on the left, full-text logo on the right. \projectname is
% still set because the cover title uses it.
\projectname{Eclipse NMBU}
\projectlogo{../../../Common/Eclipse Images/EclipseMoon}
\orglogo{../../../Common/Eclipse Images/EclipseFullText}

%
% Image paths are relative to the compiled document's folder, so every
% document using this file must sit one folder below it.
% ---------------------------------------------------------------------------

% Meeting details, used on covers and in each document's ID block.
\newcommand{\meetingdate}{30 September 2026}
\newcommand{\meetingtime}{12:00}
\newcommand{\meetinglocation}{Festsalen, NMBU}
\newcommand{\meetingline}{\textbf{Meeting:} \meetingdate, \meetingtime, \meetinglocation}

% Header: moon logo on the left, full-text logo on the right. \projectname is
% still set because the cover title uses it.
\projectname{Eclipse NMBU}
\projectlogo{../../../Common/Eclipse Images/EclipseMoon}
\orglogo{../../../Common/Eclipse Images/EclipseFullText}

%
% Image paths are relative to the compiled document's folder, so every
% document using this file must sit one folder below it.
% ---------------------------------------------------------------------------

% Meeting details, used on covers and in each document's ID block.
\newcommand{\meetingdate}{30 September 2026}
\newcommand{\meetingtime}{12:00}
\newcommand{\meetinglocation}{Festsalen, NMBU}
\newcommand{\meetingline}{\textbf{Meeting:} \meetingdate, \meetingtime, \meetinglocation}

% Header: moon logo on the left, full-text logo on the right. \projectname is
% still set because the cover title uses it.
\projectname{Eclipse NMBU}
\projectlogo{../../../Common/Eclipse Images/EclipseMoon}
\orglogo{../../../Common/Eclipse Images/EclipseFullText}

\usepackage{enumitem}

\reporttitle{Organisational Glossary}

% ---------------------------------------------------------------------------
% Cover / title page for rapport-class documents.
%
% Usage (after \documentclass{rapport} and before \begin{document}):
%   % ---------------------------------------------------------------------------
% Cover / title page for rapport-class documents.
%
% Usage (after \documentclass{rapport} and before \begin{document}):
%   % ---------------------------------------------------------------------------
% Cover / title page for rapport-class documents.
%
% Usage (after \documentclass{rapport} and before \begin{document}):
%   \input{../cover/rapport-cover}
%   \missionpatch{path/to/patch}      % optional circular mission/team patch;
%                                     % without it, the org logo is shown large
%                                     % at the top instead
%   \coverkind{Vedtekter}             % optional small line above the name
%   \reportsubtitle{Optional subtitle line}
%   \orgname{Eclipse NMBU}{Norwegian University of Life Sciences, Norway}
%   \reportlocation{Ås, Norge}
%   \date{3. oktober 2022}            % optional, defaults to \today
%   \contactinfo{Contact: post@eclipsenmbu.no}  % optional, shown at the bottom
%
% Then, as the first thing inside \begin{document}:
%   \makecoverpage
%
% The large title line shows \projectname (falling back to \reporttitle if
% \coverkind is not set), so a document typically sets both:
%   \projectname{Eclipse NMBU}
%   \coverkind{Vedtekter}
% ---------------------------------------------------------------------------

\makeatletter
\def\@missionpatch{}
\newcommand{\missionpatch}[1]{\def\@missionpatch{#1}}

\def\@coverkind{}
\newcommand{\coverkind}[1]{\def\@coverkind{#1}}

\def\@reportsubtitle{}
\newcommand{\reportsubtitle}[1]{\def\@reportsubtitle{#1}}

\def\@orgnameone{}
\def\@orgnametwo{}
\newcommand{\orgname}[2]{\def\@orgnameone{#1}\def\@orgnametwo{#2}}

\def\@reportlocation{}
\newcommand{\reportlocation}[1]{\def\@reportlocation{#1}}

\def\@contactinfo{}
\newcommand{\contactinfo}[1]{\def\@contactinfo{#1}}

\newcommand{\makecoverpage}{%
  \begingroup
  \thispagestyle{empty}
  \begin{center}
  \ifx\@missionpatch\@empty
    \includegraphics[width=10cm]{\@orglogo}\par
  \else
    \includegraphics[width=6.5cm]{\@missionpatch}\par
  \fi
  \vspace{3.4cm}
  \ifx\@coverkind\@empty
    {\sffamily\Huge\bfseries \@reporttitle \par}
  \else
    {\sffamily\Huge\bfseries \@coverkind \par}
    {\sffamily\Huge\bfseries \@projectname \par}
  \fi
  \ifx\@reportsubtitle\@empty\else
    \vspace{1.6cm}
    {\sffamily\bfseries \@reportsubtitle \par}
  \fi
  \vspace{2cm}
  \ifx\@missionpatch\@empty\else
    \includegraphics[height=1.3cm]{\@orglogo}\par
    \vspace{0.6cm}
  \fi
  \@orgnameone \\
  \@orgnametwo \par
  \vspace{0.6cm}
  \ifx\@reportlocation\@empty
    \@date
  \else
    \@reportlocation, \@date
  \fi
  \par
  \vfill
  \ifx\@contactinfo\@empty\else
    {\small \@contactinfo \par}
  \fi
  \end{center}
  \endgroup
  \newpage
  \pagenumbering{roman}
  \pagestyle{plain}
}
\makeatother

%   \missionpatch{path/to/patch}      % optional circular mission/team patch;
%                                     % without it, the org logo is shown large
%                                     % at the top instead
%   \coverkind{Vedtekter}             % optional small line above the name
%   \reportsubtitle{Optional subtitle line}
%   \orgname{Eclipse NMBU}{Norwegian University of Life Sciences, Norway}
%   \reportlocation{Ås, Norge}
%   \date{3. oktober 2022}            % optional, defaults to \today
%   \contactinfo{Contact: post@eclipsenmbu.no}  % optional, shown at the bottom
%
% Then, as the first thing inside \begin{document}:
%   \makecoverpage
%
% The large title line shows \projectname (falling back to \reporttitle if
% \coverkind is not set), so a document typically sets both:
%   \projectname{Eclipse NMBU}
%   \coverkind{Vedtekter}
% ---------------------------------------------------------------------------

\makeatletter
\def\@missionpatch{}
\newcommand{\missionpatch}[1]{\def\@missionpatch{#1}}

\def\@coverkind{}
\newcommand{\coverkind}[1]{\def\@coverkind{#1}}

\def\@reportsubtitle{}
\newcommand{\reportsubtitle}[1]{\def\@reportsubtitle{#1}}

\def\@orgnameone{}
\def\@orgnametwo{}
\newcommand{\orgname}[2]{\def\@orgnameone{#1}\def\@orgnametwo{#2}}

\def\@reportlocation{}
\newcommand{\reportlocation}[1]{\def\@reportlocation{#1}}

\def\@contactinfo{}
\newcommand{\contactinfo}[1]{\def\@contactinfo{#1}}

\newcommand{\makecoverpage}{%
  \begingroup
  \thispagestyle{empty}
  \begin{center}
  \ifx\@missionpatch\@empty
    \includegraphics[width=10cm]{\@orglogo}\par
  \else
    \includegraphics[width=6.5cm]{\@missionpatch}\par
  \fi
  \vspace{3.4cm}
  \ifx\@coverkind\@empty
    {\sffamily\Huge\bfseries \@reporttitle \par}
  \else
    {\sffamily\Huge\bfseries \@coverkind \par}
    {\sffamily\Huge\bfseries \@projectname \par}
  \fi
  \ifx\@reportsubtitle\@empty\else
    \vspace{1.6cm}
    {\sffamily\bfseries \@reportsubtitle \par}
  \fi
  \vspace{2cm}
  \ifx\@missionpatch\@empty\else
    \includegraphics[height=1.3cm]{\@orglogo}\par
    \vspace{0.6cm}
  \fi
  \@orgnameone \\
  \@orgnametwo \par
  \vspace{0.6cm}
  \ifx\@reportlocation\@empty
    \@date
  \else
    \@reportlocation, \@date
  \fi
  \par
  \vfill
  \ifx\@contactinfo\@empty\else
    {\small \@contactinfo \par}
  \fi
  \end{center}
  \endgroup
  \newpage
  \pagenumbering{roman}
  \pagestyle{plain}
}
\makeatother

%   \missionpatch{path/to/patch}      % optional circular mission/team patch;
%                                     % without it, the org logo is shown large
%                                     % at the top instead
%   \coverkind{Vedtekter}             % optional small line above the name
%   \reportsubtitle{Optional subtitle line}
%   \orgname{Eclipse NMBU}{Norwegian University of Life Sciences, Norway}
%   \reportlocation{Ås, Norge}
%   \date{3. oktober 2022}            % optional, defaults to \today
%   \contactinfo{Contact: post@eclipsenmbu.no}  % optional, shown at the bottom
%
% Then, as the first thing inside \begin{document}:
%   \makecoverpage
%
% The large title line shows \projectname (falling back to \reporttitle if
% \coverkind is not set), so a document typically sets both:
%   \projectname{Eclipse NMBU}
%   \coverkind{Vedtekter}
% ---------------------------------------------------------------------------

\makeatletter
\def\@missionpatch{}
\newcommand{\missionpatch}[1]{\def\@missionpatch{#1}}

\def\@coverkind{}
\newcommand{\coverkind}[1]{\def\@coverkind{#1}}

\def\@reportsubtitle{}
\newcommand{\reportsubtitle}[1]{\def\@reportsubtitle{#1}}

\def\@orgnameone{}
\def\@orgnametwo{}
\newcommand{\orgname}[2]{\def\@orgnameone{#1}\def\@orgnametwo{#2}}

\def\@reportlocation{}
\newcommand{\reportlocation}[1]{\def\@reportlocation{#1}}

\def\@contactinfo{}
\newcommand{\contactinfo}[1]{\def\@contactinfo{#1}}

\newcommand{\makecoverpage}{%
  \begingroup
  \thispagestyle{empty}
  \begin{center}
  \ifx\@missionpatch\@empty
    \includegraphics[width=10cm]{\@orglogo}\par
  \else
    \includegraphics[width=6.5cm]{\@missionpatch}\par
  \fi
  \vspace{3.4cm}
  \ifx\@coverkind\@empty
    {\sffamily\Huge\bfseries \@reporttitle \par}
  \else
    {\sffamily\Huge\bfseries \@coverkind \par}
    {\sffamily\Huge\bfseries \@projectname \par}
  \fi
  \ifx\@reportsubtitle\@empty\else
    \vspace{1.6cm}
    {\sffamily\bfseries \@reportsubtitle \par}
  \fi
  \vspace{2cm}
  \ifx\@missionpatch\@empty\else
    \includegraphics[height=1.3cm]{\@orglogo}\par
    \vspace{0.6cm}
  \fi
  \@orgnameone \\
  \@orgnametwo \par
  \vspace{0.6cm}
  \ifx\@reportlocation\@empty
    \@date
  \else
    \@reportlocation, \@date
  \fi
  \par
  \vfill
  \ifx\@contactinfo\@empty\else
    {\small \@contactinfo \par}
  \fi
  \end{center}
  \endgroup
  \newpage
  \pagenumbering{roman}
  \pagestyle{plain}
}
\makeatother

\coverkind{Multilingual Organisational Glossary}
\reportsubtitle{Flerspråklig organisasjonsordliste: 26H-ECL-GA-GLOSS1}
\orgname{Eclipse NMBU}{Norwegian University of Life Sciences, Norway}
\reportlocation{\meetinglocation, Ås}
\date{\meetingdate, \meetingtime}
\contactinfo{Contact: post@eclipsenmbu.no}

\newcommand{\term}[1]{\vsubsection{#1}}
\newcommand{\enDef}[1]{\noindent\textbf{EN:} #1\par}
\newcommand{\noDef}[1]{\noindent\textbf{NO:} #1\par}
\newcommand{\lastupdated}[1]{\noindent\textit{\small Last updated: #1}\par\vspace{0.4cm}}

\begin{document}

\makecoverpage
\tableofcontents
\reportmainmatter

\noindent\textbf{Document ID:} 26H-ECL-GA-GLOSS1\\
\meetingline\\
\textbf{Status:} Living reference document: update whenever a new organisational term is adopted, and update the ``Last updated'' date on any term whose translation changes.\\
\textbf{Scope:} Applies to all Eclipse NMBU governance documents, including this General Assembly's document set (26H-ECL-GA-*). Where the Old and New Bylaws use the same Norwegian term with different meaning, both are noted.

\vsection{How to use this glossary}
Every organisational document should write technical terminology as \textbf{English (Norwegian)} on first use per section, e.g. ``General Assembly (generalforsamling)''. This glossary is the authoritative source for that pairing. If a term is not listed here, propose an addition rather than inventing an ad-hoc translation.

\vsection{Organs and bodies}

\term{General Assembly (generalforsamling)}
\enDef{The organisation's supreme decision-making body, made up of its members.}
\noDef{Foreningens øverste organ, sammensatt av medlemmene.}
\lastupdated{2026-09-23}

\term{Board (styret)}
\enDef{The organisation's highest organ between General Assemblies.}
\noDef{Foreningens øverste organ mellom generalforsamlingene.}
\lastupdated{2026-09-23}

\term{Election and Control Committee (valgog kontrollkomit\'e)}
\enDef{A three-member committee elected by the General Assembly (New Bylaws only) that finds candidates, runs elections to elected offices, and oversees Bylaw compliance. Does not exist under the Old Bylaws.}
\noDef{Komit\'e p\r{a} tre medlemmer valgt av generalforsamlingen (kun Nye Vedtekter) som finner kandidater, gjennomfører valg til tillitsverv, og p\r{a}ser at vedtektene blir fulgt.}
\lastupdated{2026-09-23}

\vsection{Roles (C-suite)}

\term{CEO: Chief Executive Officer (Styreleder)}
\enDef{Overall administrative leadership and external representation. Under the New Bylaws, sole holder of Signature Right, with a duty to delegate Authority to the other Board members.}
\noDef{Overordnet administrativ ledelse og ekstern representasjon. Under Nye Vedtekter eneinnehaver av signaturrett, med plikt til \r{a} delegere fullmakt til de øvrige styremedlemmene.}
\lastupdated{2026-09-23}

\term{CTO: Chief Technology Officer / Chief Technical Officer (Teknisk ansvarlig)}
\enDef{Responsible for all technical projects, system integration and technical safety. Under the New Bylaws, also Deputy Chair (nestleder) who assumes the CEO's responsibilities when the CEO is unavailable.}
\noDef{Ansvarlig for alle tekniske prosjekt, systemintegrasjon og teknisk sikkerhet. Under Nye Vedtekter ogs\r{a} nestleder, som overtar CEOs ansvar n\r{a}r CEO er fraværende.}
\lastupdated{2026-09-23}

\term{CFO: Chief Financial Officer (Økonomiansvarlig)}
\enDef{Old Bylaws role: financial planning, bookkeeping, budget control. Replaced by CBO under the New Bylaws.}
\noDef{Rolle under Gamle Vedtekter: finansiell planlegging, bokføring, budsjettkontroll. Erstattes av CBO under Nye Vedtekter.}
\lastupdated{2026-09-23}

\term{CMO: Chief Marketing Officer (Markedsføringsansvarlig)}
\enDef{Old Bylaws role: recruitment, branding, sponsor relations, external communication. Folded into CBO under the New Bylaws.}
\noDef{Rolle under Gamle Vedtekter: rekruttering, merkevarebygging, sponsorrelasjoner, ekstern kommunikasjon. Sl\r{a}s sammen med CBO under Nye Vedtekter.}
\lastupdated{2026-09-23}

\term{CBO: Chief Business Officer (Forretningsog økonomiansvarlig)}
\enDef{New Bylaws role only. Holds full authority over the organisation's economy (undelegated from the CEO) plus sponsor relations, branding and external digital communication: effectively merging the Old Bylaws' CFO and CMO roles.}
\noDef{Kun en rolle under Nye Vedtekter. Har full myndighet over foreningens økonomi (ikke delegert fra CEO), samt sponsorrelasjoner, merkevarebygging og ekstern digital kommunikasjon: sl\r{a}r i praksis sammen CFOog CMO-rollene fra Gamle Vedtekter.}
\lastupdated{2026-09-23}

\term{CPO: Chief People Officer (Personalog medlemsansvarlig)}
\enDef{New Bylaws role only. Responsible for members and internal organisation: recruitment, onboarding, membership-tier administration, internal communication, and handling of whistleblowing cases together with the CEO.}
\noDef{Kun en rolle under Nye Vedtekter. Ansvarlig for medlemmer og intern organisasjon: rekruttering, opptak, forvaltning av medlemsniv\r{a}er, internkommunikasjon, og saksbehandling av varsler sammen med CEO.}
\lastupdated{2026-09-23}

\term{Project Lead (prosjektleder)}
\enDef{Member with overall responsibility for carrying out a project. Reports to the CTO for technical projects, otherwise to the Board member responsible for that subject area.}
\noDef{Medlem med det overordnede ansvaret for gjennomføringen av et prosjekt. Rapporterer til CTO i tekniske prosjekter, ellers til ansvarlig styremedlem.}
\lastupdated{2026-09-23}

\term{Team Lead (teamleder)}
\enDef{Member with responsibility for a team: a cross-project subject-matter group. Reports to the Project Lead of the relevant project, or to the responsible Board member.}
\noDef{Medlem med ansvar for et team, det vil si en faggruppe som g\r{a}r p\r{a} tvers av prosjekter. Rapporterer til prosjektlederen i det aktuelle prosjektet, eller til ansvarlig styremedlem.}
\lastupdated{2026-09-23}

\term{Nearest Leader (nærmeste leder)}
\enDef{The role a member reports to in the organisation's governance structure: e.g. a Team Lead for an ordinary member, a Project Lead for a Team Lead. Used to determine escalation when a Responsibility holder is unavailable.}
\noDef{Den rollen et medlem rapporterer til i foreningens styringsstruktur. Brukes for \r{a} avgjøre eskalering n\r{a}r en ansvarshaver er fraværende.}
\lastupdated{2026-09-23}

\vsection{Membership tiers}

\term{Board Member (styremedlem)}
\enDef{Holds the legal and administrative responsibility for the organisation. Highest membership tier.}
\noDef{Innehar det juridiske og administrative ansvaret for foreningen. Høyeste medlemsniv\r{a}.}
\lastupdated{2026-09-23}

\term{Core Member (kjernemedlem)}
\enDef{A member who has been assigned a genuine Responsibility Role by the Board or a Board member, e.g. as Project Lead or Team Lead.}
\noDef{Medlem som er tildelt et ansvarsverv av styret eller et styremedlem, f.eks. som prosjektleder eller teamleder.}
\lastupdated{2026-09-23}

\term{Active Member (aktivt medlem)}
\enDef{A member admitted under the ordinary admission process (New Bylaws \S5.2); has voting rights and is eligible for office. Under the Old Bylaws, an active member is instead defined by meeting a points-earning requirement.}
\noDef{Medlem tatt opp etter ordinær opptaksprosess (Nye Vedtekter \S5.2); har stemmerett og er valgbar. Under Gamle Vedtekter defineres aktivt medlem i stedet av oppfylt poengkrav.}
\lastupdated{2026-09-23}

\term{Mentor (mentor)}
\enDef{New Bylaws only. A membership tier for people, including non-members, on the organisation's mentor list.}
\noDef{Kun under Nye Vedtekter. Medlemsniv\r{a} for personer, ogs\r{a} utenfor foreningen, som st\r{a}r p\r{a} foreningens mentorliste.}
\lastupdated{2026-09-23}

\term{Inactive Member (inaktivt medlem)}
\enDef{Registered member who does not meet the requirements for active participation. Has meeting and speaking rights but no voting rights or eligibility for office.}
\noDef{Registrert medlem som ikke oppfyller kravene til aktiv deltakelse. Har møteog talerett, men ikke stemmerett eller valgbarhet.}
\lastupdated{2026-09-23}

\term{Alumni (alumni)}
\enDef{New Bylaws only. A member whose Membership Contract ended without being renewed, automatically receives Alumni status.}
\noDef{Kun under Nye Vedtekter. Et medlem hvis medlemskontrakt er avsluttet uten \r{a} bli fornyet, f\r{a}r automatisk alumnistatus.}
\lastupdated{2026-09-23}

\vsection{Governance, responsibility, delegation, signature}

\term{Responsibility (ansvar)}
\enDef{A genuine leadership obligation assigned to a member, distinct from an ordinary task. Receiving Responsibility makes a member a Core Member.}
\noDef{En reell lederforpliktelse tildelt et medlem, til forskjell fra en ordinær oppgave. \r{a} motta ansvar gjør et medlem til kjernemedlem.}
\lastupdated{2026-09-23}

\term{Responsibility Letter / Contract (ansvarsbrev / ansvarskontrakt)}
\enDef{The signed document that transfers a Responsibility from one person (``From'') to another (``To''), logged afterwards by the Board.}
\noDef{Det signerte dokumentet som overfører et ansvar fra en person (``Fra'') til en annen (``Til''), som deretter logges av styret.}
\lastupdated{2026-09-23}

\term{Authority / Mandate (fullmakt)}
\enDef{The right to make decisions and bind the organisation on the CEO's behalf, within an assigned area of responsibility.}
\noDef{Retten til \r{a} treffe beslutninger og forplikte foreningen p\r{a} vegne av CEO innenfor et angitt ansvarsomr\r{a}de.}
\lastupdated{2026-09-23}

\term{Delegation (delegasjon)}
\enDef{Transfer of Authority from the CEO to another Board member (New Bylaws \S14.2), or onward from a Board member to a Project Lead or Team Lead for defined tasks (\S14.6).}
\noDef{Overføring av fullmakt fra CEO til et annet styremedlem (Nye Vedtekter \S14.2), eller videre fra et styremedlem til en prosjektleder eller teamleder for definerte oppgaver (\S14.6).}
\lastupdated{2026-09-23}

\term{Signature Right (signaturrett)}
\enDef{The right to sign binding agreements on the organisation's behalf. Old Bylaws: held jointly by the Board. New Bylaws: held solely by the CEO, who may delegate it directly to other Board members.}
\noDef{Retten til \r{a} signere bindende avtaler p\r{a} vegne av foreningen. Gamle Vedtekter: innehas i fellesskap av styret. Nye Vedtekter: innehas alene av CEO, som kan delegere den direkte til andre styremedlemmer.}
\lastupdated{2026-09-23}

\vsection{Meetings and voting}

\term{Meeting Chair (ordstyrer)}
\enDef{Chairs the General Assembly, manages the speakers' list and enforces meeting procedure.}
\noDef{Leder generalforsamlingen, fører taleliste og h\r{a}ndhever møtereglene.}
\lastupdated{2026-09-23}

\term{Secretary / Minutes Taker (referent)}
\enDef{Records the Protocol of the General Assembly.}
\noDef{Fører protokoll fra generalforsamlingen.}
\lastupdated{2026-09-23}

\term{Vote Counters (tellekorps)}
\enDef{Two neutral members who count votes not otherwise counted by the Election and Control Committee, and record the number of eligible voters present.}
\noDef{To nøytrale medlemmer som teller stemmer som ikke telles av valgog kontrollkomiteen, og registrerer antall fremmøtte stemmeberettigede.}
\lastupdated{2026-09-23}

\term{Protocol Signers (protokollunderskrivere)}
\enDef{Two members elected to check and sign the Protocol as an accurate record of the General Assembly.}
\noDef{To medlemmer valgt til \r{a} kontrollere og signere protokollen som en korrekt fremstilling av generalforsamlingen.}
\lastupdated{2026-09-23}

\term{Case (sak)}
\enDef{An item on the Agenda to be presented, discussed and resolved by the General Assembly.}
\noDef{Et punkt p\r{a} sakslisten som skal presenteres, diskuteres og avgjøres av generalforsamlingen.}
\lastupdated{2026-09-23}

\term{Case Paper (sakspapir)}
\enDef{The written document prepared for a Case, following the Case Template.}
\noDef{Det skriftlige dokumentet utarbeidet for en sak, etter saksmalen.}
\lastupdated{2026-09-23}

\term{Right to Propose (forslagsrett)}
\enDef{The right to submit Cases and Amendments to the General Assembly. Held by all members with voting rights, and by Mentors.}
\noDef{Retten til \r{a} melde inn saker og fremme endringsforslag til generalforsamlingen. Innehas av alle medlemmer med stemmerett, samt mentorer.}
\lastupdated{2026-09-23}

\term{Amendment (endringsforslag)}
\enDef{A proposal to change a Case already on the Agenda, submitted in writing before the deadline.}
\noDef{Forslag om endring av en sak som allerede st\r{a}r p\r{a} sakslisten, fremmet skriftlig før fristen.}
\lastupdated{2026-09-23}

\term{Floor Amendment (benkeforslag)}
\enDef{At the General Assembly itself: (a) in the ordinary sense, a new candidate nominated from the floor during a Case's presentation and discussion, per New Bylaws \S11.3; the term is not used in these documents for last-minute changes to Case text, which are not permitted after the Amendment deadline.}
\noDef{P\r{a} selve generalforsamlingen: (a) i vanlig forstand, en ny kandidat nominert fra salen under presentasjon og diskusjon av en sak, jf. Nye Vedtekter \S11.3; begrepet brukes ikke i disse dokumentene om sist-liten endringer i saksteksten, som ikke er tillatt etter fristen for endringsforslag.}
\lastupdated{2026-09-23}

\term{Advance Nomination (forhåndsbenking)}
\enDef{Nomination of a candidate before the General Assembly, in response to the Election and Control Committee's (or, under the Old Bylaws, the Board's) call for candidates.}
\noDef{Nominasjon av en kandidat før generalforsamlingen, som svar p\r{a} valgog kontrollkomiteens (eller, under Gamle Vedtekter, styrets) oppfordring til forh\r{a}ndsbenking.}
\lastupdated{2026-09-23}

\term{Eligible Voter (stemmeberettiget)}
\enDef{A member with voting rights at the General Assembly: Board Members, Core Members and Active Members.}
\noDef{Et medlem med stemmerett p\r{a} generalforsamlingen: styremedlemmer, kjernemedlemmer og aktive medlemmer.}
\lastupdated{2026-09-23}

\term{Voting Order (voteringsorden)}
\enDef{The sequence in which proposals on a Case are put to a vote, structured to require as few individual votes as possible while resolving every proposal unambiguously.}
\noDef{Rekkefølgen forslagene til en sak stemmes over i, strukturert for \r{a} kreve s\r{a} f\r{a} enkeltavstemninger som mulig, samtidig som hvert forslag avgjøres entydig.}
\lastupdated{2026-09-23}

\term{Constitution of the Meeting (konstituering)}
\enDef{The opening case of the General Assembly: election of Meeting Chair, Secretary, Vote Counters and Protocol Signers, approval of the Notice and Agenda, and verification of eligible voters present.}
\noDef{\r{A}pningssaken p\r{a} generalforsamlingen: valg av ordstyrer, referent, tellekorps og protokollunderskrivere, godkjenning av innkalling og saksliste, og verifisering av fremmøtte stemmeberettigede.}
\lastupdated{2026-09-23}

\vsection{Other}

\term{Bylaws (vedtekter)}
\enDef{The organisation's founding governance document. See ``Old Bylaws'' and ``New Bylaws''.}
\noDef{Foreningens grunnleggende styringsdokument. Se ``Gamle vedtekter'' og ``Nye vedtekter''.}
\lastupdated{2026-09-23}

\term{Old Bylaws (gamle vedtekter)}
\enDef{The Bylaws in force immediately before this General Assembly, in 13 sections (\S\S1--13).}
\noDef{Vedtektene som gjelder umiddelbart før denne generalforsamlingen, i 13 paragrafer (\S\S1--13).}
\lastupdated{2026-09-23}

\term{New Bylaws (nye vedtekter)}
\enDef{The proposed restructured Bylaws, in 18 sections across five parts (\S\S1--18), submitted for adoption at this General Assembly.}
\noDef{De foresl\r{a}tte omstrukturerte vedtektene, i 18 paragrafer fordelt p\r{a} fem deler (\S\S1--18), fremmet for vedtak p\r{a} denne generalforsamlingen.}
\lastupdated{2026-09-23}

\term{Document ID (dokument-ID)}
\enDef{The stable identifier assigned to every document in this set, format \texttt{YYH/V-ECL-GA-TYPEn} (e.g.\ \texttt{26H-ECL-GA-CASE0}), allowing documents to cross-reference each other regardless of final Case numbering. See the ID and Abbreviations Reference (26H-ECL-GA-REF1) for the full scheme.}
\noDef{Den stabile identifikatoren tildelt hvert dokument i dette settet, p\r{a} formatet \texttt{YYH/V-ECL-GA-TYPEn} (f.eks.\ \texttt{26H-ECL-GA-CASE0}), som gjør det mulig for dokumenter \r{a} referere til hverandre uavhengig av endelig saksnummerering. Se ID- og forkortelsesreferansen (26H-ECL-GA-REF1) for hele systemet.}
\lastupdated{2026-09-23}

\end{document}
